\subsection{Degree of a covering}

Let B be a topological space, let E be a covering of B, and denote its projection by p. Let C be the set of cardinals \(\operatorname{Card}(\overline{p}^{1}(b))\), where b ranges over B. The map \(b \mapsto \operatorname{Card}(\overline{p}^{1}(b))\) is a locally constant map from B to C. One says that the covering E has a degree if B is nonempty and if the map \(b \mapsto \operatorname{Card}(\overline{p}^{1}(b))\) is constant. The common value of the cardinals \(\operatorname{Card}(\overline{p}^{1}(b))\), for \(b \in B\), is then called the degree of the covering E and is denoted \(\deg(E, p)\), or simply \([E:B]\) if there can be no ambiguity about the map p.

If B is not empty, the trivial covering with base B and fiber type F has a degree equal to Card(F).

If B is connected, the function \(b \mapsto \text{Card}(\overline{p}^{1}(b))\) is constant. Consequently:

PROPOSITION 4. --- Every covering of a nonempty connected space has a degree.

Let G be a topological space, let \((\mathrm{F}, g)\) be a covering of G and let \((\mathrm{E}, f)\) be a covering of F; suppose that these coverings have a degree. If the G-space \((\mathrm{E}, g \circ f)\) is a covering, then it has a degree and one has \(\deg(\mathrm{E}, g \circ f) = \deg(\mathrm{F}, g) \deg(\mathrm{E}, f)\), which can also be written

\[
[ \mathrm{E}: \mathrm{G} ] = [ \mathrm{E}: \mathrm{F} ] [ \mathrm{F}: \mathrm{G} ].
\]

Indeed, if z is a point of G, all the fibers of the map

\[
f _ {g (z)} ^ {- 1}: \stackrel {{- 1}} {{f}} (\stackrel {{- 1}} {{g}} (z)) \to \stackrel {{- 1}} {{g}} (z)
\]

have cardinality \(\deg(\mathrm{E}, f)\), and \(\overline{g}^1(z)\) has cardinality \(\deg(\mathrm{F}, g)\). The assertion therefore follows from the shepherd's principle (E, III, p. 41, prop. 9).

Let B and B' be topological spaces, let (E, p) be a covering of B and let (E', p') be a covering of B'. Suppose that these coverings have a degree; then the covering (E × E', p × p') of B × B' (I, p. 71, prop. 2) has a degree and one has:

\[
\deg (\mathrm{E} \times \mathrm{E} ^ {\prime}, p \times p ^ {\prime}) = \deg (\mathrm{E}, p) \deg (\mathrm{E} ^ {\prime}, p ^ {\prime}).
\]

Indeed, for every pair \((b,b^{\prime})\in\mathrm{B}\times\mathrm{B}^{\prime}\), the fiber \((\mathrm{E}\times\mathrm{E}^{\prime})_{(b,b^{\prime})}\) is the product \(E_{b}\times E_{b^{\prime}}^{\prime}\) .

Let B and B' be topological spaces, let (E, p) be a covering of B, let (E', p') be a covering of B', and let

\[
\begin{array}{c} \mathrm{E} ^ {\prime} \xrightarrow {f ^ {\prime}} \mathrm{E} \\ \Biggl \downarrow p ^ {\prime} \\ \mathrm{B} ^ {\prime} \xrightarrow {f} \mathrm{B} \end{array}
\]

a Cartesian square. If B′ is nonempty and if the covering (E, p) has a degree, then the covering (E′, p′) has a degree, equal to that of E, by I, p. 10, cor. of prop. 4.

